\section{General lessons}

The common lesson is not that shape is irrelevant. The spectra and conductance
curves do change with geometry, and the direct calculations are useful for
controlled numerical physics. The lesson is that a shape-dependent observable is
not automatically a shape-design rule. In this project, the strongest
interpretation is often removed by a baseline that is simpler than the proposed
objective.

The first layer is confinement scale. Because $E_{\mathrm{kin}}$ near the
bottom of the band follows a continuum-like quadratic dispersion, many low
levels and gaps inherit a leading size dependence. Any residual objective must
show that it is not merely reporting this scale. The second layer is geometry
compression: area, perimeter, aspect ratio, and site count can all absorb
substantial variation before one reaches a shape-exponent-specific claim. The
third layer is symmetry and degeneracy. Q, S, and magnetic ranking ideas all
encounter this issue in different forms; isotropy and degeneracy lifting can be
physically real while still being too simple to count as inverse design.

The fourth layer is boundary realization. A smooth superellipse becomes a
finite set of lattice sites, so boundary pixelation, area error, and sublattice
imbalance can produce finite-size structure. These effects are not nuisances to
hide; they are part of the model and must be included in the baseline ladder.
The fifth layer is reference-solver uncertainty. The FD-reference sequence shows
that passing a circular sanity check is not enough for a shape-sensitive
continuum residual. A reference can be useful for calibration and still be too
uncertain for the intended residual analysis.

Transport adds one more layer: contacts. Conductance is not only a property of
the isolated dot shape. It depends on the coupling between the finite dot and
the leads, on lead width, and on mode structure. The preflight therefore does
what a preflight should do: it blocks a large shape-transport campaign until
the lead-dot coupling baseline is controlled.

These lessons also clarify the role of ML. A surrogate can be valuable for
compression, interpolation, and candidate generation, especially when direct
Kwant calculations are more expensive than model evaluation. But in the present
dataset, a flexible model is not automatically more scientific than a Ridge
baseline tied to confinement descriptors. ML should enter after the physical
question survives baseline subtraction, not as a substitute for that survival.

The resulting baseline hierarchy is:
\begin{enumerate}
    \item check size, area, perimeter, and aspect-ratio scaling;
    \item check isotropy, symmetry, and degeneracy;
    \item check lattice boundary realization, pixelation, and effective radius;
    \item check reference-solver uncertainty before using residuals;
    \item check contact geometry before interpreting transport;
    \item only then use ML or inverse screening for positive claims.
\end{enumerate}
Applied to the existing outputs, this hierarchy supports a reproducible
negative benchmark rather than a new device-design rule.
